\chapter{Numerical Root Finding}

\section{Why Approximate a Root?}\index{root!numerical}
Your calculator has a $\sqrt{\ }$ button, but nobody stored the answer
for every number inside it. When you press it, the chip makes a guess, checks
how wrong it is, and improves the guess, a few times, in a fraction of a
millisecond. This chapter shows you two ways to do that.

A \newterm{root} of a function $f$ is a value $x$ where $f(x) = 0$. For
$f(x) = x^2 - 2$, the root is $\sqrt{2}$. Many equations have no formula for their
roots, such as $\cos x = x$ or $e^x = 3x$. For those we use numerical methods:
we start with a guess and improve it until it is accurate enough.

\section{The Bisection Method}\index{bisection method}
Bisection uses the Intermediate Value Theorem (see the limits chapter). If $f$ is
continuous and $f(a)$ and $f(b)$ have opposite signs, then $f$ has a root between
$a$ and $b$. Bisection checks the midpoint and keeps the half that still has a sign change.

For $f(x) = x^2 - 2$ on $[1, 2]$, we have $f(1) = -1$ and $f(2) = 2$. The
midpoint is $1.5$, and $f(1.5) = 0.25 > 0$. The sign changes between $1$ and
$1.5$, so the new interval is $[1, 1.5]$. Each step halves the interval. To
find $\sqrt{2}$ to within $10^{-6}$, bisection needs 19 steps.

\section{Newton's Method}\index{Newton's method}
Newton's method uses the tangent line. Start at $x_0$. The tangent line at
$x_0$ crosses the $x$-axis at
$$x_{n+1} = x_n - \frac{f(x_n)}{f'(x_n)}$$
For $f(x) = x^2 - 2$, this is $x_{n+1} = x_n - \frac{x_n^2 - 2}{2x_n}$. Starting
at $x_0 = 1$:
$$x_1 = 1.5, \qquad x_2 = 1.41\overline{6}, \qquad x_3 \approx 1.414216$$
Starting from the same point, the iterates stop changing by less than $10^{-12}$
after 6 steps. It works quickly when the starting guess is close, but it can fail if
$f'(x_n)$ is zero or the guess is far from the root.

\textbf{Example}: Solve $\cos x = x$.

\textbf{Solution}: Rewrite it as $f(x) = \cos x - x = 0$, so
$f'(x) = -\sin x - 1$. Starting at $x_0 = 1$:
$$x_1 \approx 0.750364, \qquad x_2 \approx 0.739113, \qquad x_3 \approx 0.739085, \qquad x_4 \approx 0.739085$$
Three steps give six correct decimal places. No algebra trick would have found
this number.

Newton's method can also get stuck. For $f(x) = x^3 - 2x + 2$ starting at
$x_0 = 0$, you get $x_1 = 1$, then $x_2 = 0$, then $x_3 = 1$, forever. The
tangent at each point aims straight at the other one. Always check that your
iterates are actually settling down.

\section{Code}\index{bisection method!code}
\begin{minted}{python}
def bisection(f, a, b, tol=1e-6):
    if f(a) * f(b) > 0:
        raise ValueError("f(a) and f(b) must have opposite signs")
    while (b - a) / 2 > tol:
        mid = (a + b) / 2
        if f(a) * f(mid) <= 0:
            b = mid
        else:
            a = mid
    return (a + b) / 2

def newton(f, df, x, tol=1e-10, max_steps=50):
    for _ in range(max_steps):
        x_next = x - f(x) / df(x)
        if abs(x_next - x) < tol:
            return x_next
        x = x_next
    raise RuntimeError("Newton's method did not converge")

f = lambda x: x**2 - 2
df = lambda x: 2 * x
print(bisection(f, 1, 2))  # 1.4142141342163086
print(newton(f, df, 1.0))  # 1.4142135623730951
\end{minted}

\section{Comparing the Methods}
Bisection always converges when it starts with a sign change, but it is slow. It
gains about one correct binary digit per step. Newton's method usually converges
much faster, but it needs the derivative, and it may diverge from a poor starting point.

\begin{Exercise}[label=rootfind1, title=One bisection step]
For $f(x) = x^2 - 2$ on $[1, 2]$, find the new interval after one bisection step.
\end{Exercise}
\begin{Answer}[ref=rootfind1]
The midpoint is $1.5$ and $f(1.5) = 0.25 > 0$, so the sign change is on $[1, 1.5]$.
\end{Answer}

\begin{Exercise}[label=rootfind2, title=One Newton step]
Starting at $x_0 = 2$, find $x_1$ and $x_2$ for $f(x) = x^2 - 2$.
\end{Exercise}
\begin{Answer}[ref=rootfind2]
$x_1 = 2 - \frac{2}{4} = 1.5$, and $x_2 = 1.5 - \frac{0.25}{3} \approx 1.4167$.
\end{Answer}

\begin{Exercise}[label=rootfind3, title=No sign change]
Why does \pyfunction{bisection} raise an error for $f(x) = x^2 + 1$ on $[-1, 1]$?
\end{Exercise}
\begin{Answer}[ref=rootfind3]
$f(-1) = 2$ and $f(1) = 2$ have the same sign, so the Intermediate Value Theorem
does not guarantee a root in the interval. (The function has no real roots at all.)
\end{Answer}

\begin{Exercise}[label=rootfind4, title=How many halvings?]
You run bisection on an interval of width 1. How many steps until the interval
is narrower than $0.001$?
\end{Exercise}
\begin{Answer}[ref=rootfind4]
After $n$ steps the width is $\frac{1}{2^n}$. Since $2^{10} = 1024 > 1000$, you
need $10$ steps.
\end{Answer}

